%==============================================================================
\chapter{Functional description of needs}
%==============================================================================

\section{Entity management (CRUD)}
The system must ensure the creation, viewing, modification and deletion of all
business entities: farmers, farms, sites, hives, bodies/supers, frames, agents,
visits and reports, in compliance with the management rules stated in
chapter~3.

\section{Planning and monitoring of visits}
\begin{itemize}
  \item Visits are \textbf{regular or irregular}, for various reasons:
        health inspection, populating with a new swarm, requeening,
        food replenishment, medication prescription, swarm splitting,
        adding a frame and/or an extension level, collecting honey frames,
        production assessment, assessment of the swarm's population and volume.
  \item Each agent follows a \textbf{visit schedule approved} by the supervising agent of the hive concerned.
\end{itemize}

\subsection{Visit report}
Each visit gives rise to a report containing: \textbf{date, time, duration,
reason, findings, planned actions, actions carried out, recommendations}, a
qualitative assessment of the swarm's population / volume, its health status and
its productivity level on a \textbf{scale of 1 to 3}.

\section{Dashboards}

\subsection{Calendar / agents × hives matrix}
A matrix cross-referencing \textbf{agents and hives}, broken down by month, week,
season and year, with \textbf{filters} and groupings by site, by agent and by
responsible agent. Clicking a cell representing a planned visit displays a
\textbf{pop-up} summarising what is planned (date, time, reason, actions) and, if
the visit has already been carried out, what was done together with the actual
date and time.

\subsection{Production dashboard (farmer / owner)}
Highlights the \textbf{site × hive} pairs whose production (weight) has exceeded
a \textbf{configurable threshold}, indicating either an imminent honey harvest or
an extension and/or swarm-splitting operation to encourage colony expansion.

\subsection{Alerts dashboard (manager)}
Flags the \textbf{site × hive} pairs whose production or population is dropping
unusually beyond a configurable threshold, triggering an \textbf{alert} and an
imminent inspection.

\subsection{Summary dashboard (owner)}
Charts, curves and histograms representing:
\begin{itemize}
  \item the quantity produced per farm and per site
  \item the number of operations per farm, per site and per hive
  \item the \textbf{return on investment} (cost / production ratio) to decide
        the fate of each site or hive (closure, relocation).
\end{itemize}

\section{Performance indicators (automated part: preparation)}
The measurements are \textbf{time-stamped and geolocated}, expressed in
\textbf{configurable} units (internationalisation of unit systems). The
measurements of a single indicator do not necessarily share the same unit, as
the sensors are heterogeneous. Indicators concerned:
\begin{itemize}
  \item weight of the frame, the base and the extension level
  \item temperature and humidity, inside and outside the hive
  \item count of bee entry and exit movements
  \item noise / buzzing level, inside and outside
  \item light level, inside and outside
  \item wind speed and direction outside
  \item hive opening status (operation, storm, theft, animal attack, etc.).
\end{itemize}

\subsection{Predictive analysis and connected monitoring}
Connected monitoring solutions (precision beekeeping) such as Arnia,
BroodMinder or BeeHero show that advanced exploitation of the same indicators
makes it possible to anticipate colony events. The data model must remain open
to these processes, planned for the automated part.
\begin{itemize}
  \item acoustic detection of swarming 1 to 5 days in advance through analysis of the colony's sound signature
  \item confirmation of the queen's presence from the buzzing
  \item detection of opening, shock or tipping of the hive by accelerometer (theft, fall, animal attack), in addition to the opening indicator
  \item hive orientation by electronic compass
  \item continuous weighing by scale to track the nectar flow and production dynamics
  \item cloud-based analysis by artificial intelligence of acoustic signatures to correlate sound and hive activity.
\end{itemize}

\subsection{Measurement ingestion protocol}
The automated part remains out of the development scope, but the ingestion
interface is specified now so that the data model and the architecture
anticipate it.
\begin{itemize}
  \item \textbf{Transport}: a field gateway transmits the measurements to the
        server, chosen according to connectivity, either by \textbf{REST over
        HTTPS} (\texttt{POST /api/mesures}) or by \textbf{MQTT} (publishing on
        the topic \texttt{zumm/\{rucheId\}/\{indicateur\}}).
  \item \textbf{Format}: \textbf{JSON}, a measurement carrying the hive
        identifier, the indicator type, the value, the unit, the timestamp and
        the geolocation. A batch upload is an array of measurements.
  \item \textbf{Frequency}: configurable per indicator (for example weight
        every 15 min, temperature every 5 min, acoustics sampled
        continuously), in order to control the data volume.
  \item \textbf{Robustness}: \textbf{asynchronous} ingestion via a queue, with
        \textbf{idempotency} on the key (hive, indicator, timestamp) to replay
        without duplicates. In a remote area, local buffering on the gateway
        then deferred synchronisation.
  \item \textbf{Security}: authenticated gateway (key or certificate), TLS
        channel. Rejection of values outside the physically plausible range
        (safeguard).
\end{itemize}
\begin{verbatim}
POST /api/mesures        Content-Type: application/json
{
  "rucheId": 42, "typeIndicateur": "TEMPERATURE_INTERNE",
  "valeur": 34.8, "unite": "degC",
  "horodatage": "2026-07-09T14:05:00Z",
  "latitude": 36.806, "longitude": 10.181
}
\end{verbatim}

\subsection{Default thresholds and alert logic}
The thresholds are stored in the \texttt{SEUIL} table (configurable, chapter~5).
The values below are \textbf{reference default values}, anchored in the
beekeeping protocol (chapter~12).
\begin{xltabular}{\textwidth}{>{\bfseries}p{4.4cm} X p{4.4cm}}
\toprule
\textbf{Indicator} & \textbf{Default threshold} & \textbf{Alert triggered} \\
\midrule
\endhead
Internal temperature (brood) & \textless\,32\,°C or \textgreater\,36\,°C & Thermal risk to the brood \\
Internal humidity & \textgreater\,80\,\% & Excess humidity, health risk \\
Weight: super gain & \textgreater\ production threshold & Imminent honey harvest \\
Weight: sudden drop & \textgreater\,1.5\,kg in less than 1\,h & Swarming, theft, fall or attack \\
Production or population & Relative drop \textgreater\,20\,\% vs previous period & Imminent health inspection \\
Entry/exit activity & $\approx$\,0 on a favourable day & Weakened colony or missing queen \\
Opening status & Open outside a planned visit & Intrusion or disturbance \\
Acoustic signature & Swarming pattern detected & Swarming pre-alert (1 to 5 days) \\
\bottomrule
\end{xltabular}
\begin{encadre}
\textbf{Evaluation rule:} at each measurement, the value is converted to the
reference unit (conversion formula, §5.3), then compared with the configurable
threshold. To limit \emph{false alerts}, an alert is only raised if the breach
\textbf{persists over $N$ consecutive measurements} (debounce / hysteresis).
Any alert remains a \textbf{decision-support aid} confirmed by the visit report:
the agent keeps the final validation.
\end{encadre}

\section{Summary grid of functions}
\begin{xltabular}{\textwidth}{>{\bfseries}p{3.4cm} X C{2.2cm}}
\toprule
\textbf{Function} & \textbf{Description} & \textbf{Priority} \\
\midrule
\endhead
Entity CRUD        & Complete management of the business model                & High \\
Visit scheduling   & Planning + approval by the supervisor                    & High \\
Visit reports      & Structured recording of findings and actions             & High \\
Agents×hives calendar & Filterable matrix with summary pop-up                 & High \\
Production alerts  & Detection of threshold breaches                          & High \\
Health alerts      & Detection of abnormal population/production drops         & High \\
Summary \& ROI     & Cost/production charts, decision support                 & Medium \\
Sensor indicators  & Data model ready for automation                          & Medium \\
\midrule
\multicolumn{3}{l}{\textit{\color{mielfonce} Defect-prevention requirements (drawn from the limits of existing solutions)}} \\
\midrule
Autonomous deployment & J2EE installation with no mandatory subscription or third-party service & High \\
Offline mode & Data entry without connection and deferred synchronisation & High \\
Simple authentication & Google OAuth and fast field-entry flow & High \\
Configurable modularity & Enabling modules per profile via ConfigZumm.ini & Medium \\
Data portability & CSV and TXT export, web-service API and backup, with no lock-in & High \\
Contextual help & Simple, consistent interface, user guidance & Medium \\
Web and mobile parity & Same functions on all devices & Medium \\
Internationalisation & FR, EN and AR XML file, extensible to other languages & High \\
Hardware openness & Heterogeneous sensors and configurable units & Medium \\
Asynchronous ingestion & Deferred time-stamped measurements, autonomous manual part & Medium \\
Human validation & Configurable thresholds and confirmation by the visit report & High \\
Exchange security & SSL and X.509 certificates, authenticated access & High \\
Decision support & Indicative alerts, the agent keeps the final validation & Medium \\
\bottomrule
\end{xltabular}

\section{Additional features inspired by market solutions}
In order to strengthen usability and value in use, the system draws on the best
practices of reference beekeeping applications such as HiveTracks. These
features extend the needs already described without replacing them.
\begin{xltabular}{\textwidth}{>{\bfseries}p{4.2cm} X C{2.1cm}}
\toprule
\textbf{Function} & \textbf{Description} & \textbf{Priority} \\
\midrule
\endhead
Inspection photos & Attach one or more photos to each visit report for a visual follow-up of the colony and the frames. & High \\
Assisted voice entry & Dictate the visit report, the transcription automatically feeding the report fields. & Low \\
Weather context & Display the local weekly weather from the site location to inform operation decisions. & Medium \\
Map and foraging radii & Position the hives on a map and draw foraging radii (for example 1, 2 and 3 km, configurable units) to visualise the environment accessible to the bees. & Medium \\
Task list and reminders & Manage a task list and a reminder calendar (feeding, inspection, queen status) so as not to miss any deadline. & High \\
Queen tracking & Record the history of the queen's status (presence, age, replacement, marking) across visits. & High \\
Harvest and QR code & Record honey harvests and generate a QR code per batch presenting the tasting notes, the provenance and the photos, for traceability and valorisation. & Medium \\
Multi-device access & Offer web access and an experience suited to mobile terminals for field entry. & Medium \\
\bottomrule
\end{xltabular}
\begin{encadre}
\textbf{Consistency with the requirements:} the foraging radii and the weather
reuse the geolocation of the sites and the configurable-units mechanism. The
reminders rely on the existing visit schedule, and the queen tracking enriches
the visit report already defined.
\end{encadre}

\section{Advanced management functions (inspiration from Apiary Book and BeePlus)}
The beekeeping management applications Apiary Book and BeePlus emphasise
organisation and traceability functions that complement the system's scope.
\begin{xltabular}{\textwidth}{>{\bfseries}p{4.2cm} X C{2.1cm}}
\toprule
\textbf{Function} & \textbf{Description} & \textbf{Priority} \\
\midrule
\endhead
Offline mode and synchronisation & Enter field data without a connection and synchronise it automatically when the network returns. & High \\
Treatment tracking & Keep the history of the medications and health treatments applied to each hive. & High \\
Equipment and inventory management & Track the equipment (bodies, supers, frames, hives) and its allocation to sites. & Medium \\
Financial tracking & Record costs and revenues per farm, site and hive, in support of the return-on-investment calculation. & Medium \\
Data export & Export the records in CSV and TXT formats for external analysis. & Medium \\
Disease detection & Help identify diseases from the visit findings. & Low \\
Data backup & Back up and restore the data in the cloud. & Medium \\
\bottomrule
\end{xltabular}
\begin{encadre}
\textbf{Consistency with the requirements:} the financial tracking feeds the
return-on-investment dashboard, the CSV export extends the web-service API, and
the treatment tracking enriches the medication-prescription visit reason.
\end{encadre}

\section{Known limits of existing solutions and requirements to avoid them}
The analysis of user feedback on market solutions (HiveTracks, Apiary Book,
BeePlus) and on connected monitoring systems (Arnia, BroodMinder, BeeHero)
brings out recurring limits. The following table sets each limit against the
requirement adopted to avoid it in Zümm.
\begin{xltabular}{\textwidth}{>{\bfseries}p{5.2cm} X}
\toprule
\textbf{Observed limit} & \textbf{Requirement adopted to avoid it} \\
\midrule
\endhead
Paid subscription and paywall & Self-deployable J2EE system, with no mandatory subscription; all management functions remain accessible. \\
Limited offline operation & Robust offline mode with deferred synchronisation; field entry does not require a permanent connection. \\
Forced account creation and friction & Simple authentication via Google OAuth and a fast entry flow; usability comes first. \\
Overload of useless features & Modules enabled per profile and via the ConfigZumm.ini configuration file, thanks to the loosely coupled architecture. \\
Dependency and data loss & Data controlled by the user, CSV and TXT export, web-service API and backup, with no proprietary lock-in. \\
Dense interface and learning curve & Simple, consistent and user-friendly interface with contextual help, in line with the usability requirements. \\
Uneven functions across platforms & Functional parity between web and mobile access thanks to the multi-tier architecture. \\
English-only solutions & Internationalisation through an XML file, at least FR, EN and AR, open to other languages. \\
High cost of sensor hardware & Openness to heterogeneous and low-cost sensors with configurable units, with no hardware lock-in. \\
Autonomy and connectivity in a remote area & Asynchronous ingestion of time-stamped measurements; the manual part works without sensors. \\
False alerts and sensor reliability & Configurable thresholds and human validation through the visit report; alerts remain a decision-support aid. \\
Data confidentiality and security & SSL encryption and X.509 certificates, authenticated access. \\
Over-reliance on automation & Positioned as decision support; the beekeeping agent keeps the final validation. \\
\bottomrule
\end{xltabular}
